\section{Qualitative Inhaltsanalyse der Abschlussinterviews}
\label{sec:BB-08_qualitative-analysis}

Dieses Kapitel dokumentiert die qualitative Auswertung der leitfadengestuetzten Interviews auf Basis der qualitativen Inhaltsanalyse nach Mayring. Ziel ist die transparente Darstellung der Kategorienbildung, Codierlogik sowie der zentralen inhaltlichen Muster, die den Ergebnissen in Kapitel~5 zugrunde liegen.

% -----------------------------------------------------
\subsection{Analytisches Vorgehen}

Die Interviewdaten wurden vollstaendig transkribiert bzw. in strukturierter Stichpunktform dokumentiert. Die Auswertung erfolgte in einem deduktiv-induktiven Vorgehen:

\begin{itemize}
  \item Deduktive Hauptkategorien wurden aus den Forschungsfragen und dem Interviewleitfaden abgeleitet (RQ-1 bis RQ-3).
  \item Innerhalb dieser Hauptkategorien wurden induktiv Subkategorien aus dem Datenmaterial entwickelt.
  \item Alle Aussagen wurden fallbezogen codiert und anschliessend kategorienweise aggregiert.
  \item Mehrfachnennungen pro Teilnehmendem waren moeglich, sofern unterschiedliche inhaltliche Aspekte angesprochen wurden.
\end{itemize}

Die Codierung wurde iterativ ueberarbeitet, bis eine stabile und trennscharfe Kategorienstruktur erreicht war.

% -----------------------------------------------------
\subsection{Uebersicht der Kategorienstruktur}

\begin{table}[H]
\centering
\caption{Kategorienstruktur der qualitativen Auswertung}
\label{tab:AA-category-overview}
\begin{tabular}{p{4cm} p{8cm}}
\toprule
Hauptkategorie & Subkategorien \\
\midrule
RQ-1: Gebrauchstauglichkeit und Informationsarchitektur &
Struktur und Gliederung; Navigation und Auffindbarkeit; Visuelle Orientierung; Mentale Modelle; Priorisierung von Inhalten \\
RQ-2: Rollenpassung und UX-Strategien &
Rollenabdeckung; Einstiegsseiten; Rollendifferenzierung; Wahrgenommene Nutzerzentrierung; Generische vs. spezifische Gestaltung \\
RQ-3: Akzeptanz und organisationale Einbettung &
Nutzungsintention; Akzeptanztreiber; Akzeptanzbarrieren; Kommunikationsbedarf; Governance und Verantwortlichkeiten \\
\bottomrule
\end{tabular}
\end{table}

% -----------------------------------------------------
\subsection{RQ-1: Gebrauchstauglichkeit und Informationsarchitektur}

\textbf{Struktur und Gliederung} (8/10 TN) \\
Klare thematische Buendelung, logisch nachvollziehbare Seitenstruktur und reduzierte Komplexitaet wurden mehrfach als zentraler Verbesserungsfaktor genannt.

\textbf{Navigation und Auffindbarkeit} (7/10 TN) \\
Kuertzere Navigationswege und geringere Klickzahlen wurden als direkte Effizienzgewinne wahrgenommen.

\textbf{Visuelle Orientierung} (7/10 TN) \\
Reduziertes Design, klare Typografie und visuelle Hierarchien unterstuetzten die Orientierung.

\textbf{Mentale Modelle} (7/10 TN positiv; 7/10 TN kritisch zu Ansprechpartnern) \\
Die ueberwiegende Struktur entsprach den Erwartungen der Nutzenden. Gleichzeitig wurde die Platzierung von Ansprechpartnern als inkonsistent wahrgenommen.

\textbf{Priorisierung von Inhalten} (4/10 TN) \\
Nachrichten wurden teilweise als zu dominant bewertet; operative Inhalte sollten staerker hervorgehoben werden.

% -----------------------------------------------------
\subsection{RQ-2: Rollenpassung und UX-Strategien}

\textbf{Rollenabdeckung} (9/10 TN) \\
Alle relevanten Inhalte fuer die jeweilige Rolle wurden als vorhanden bewertet.

\textbf{Einstiegsseiten} (6/10 TN) \\
Uebersichtsseiten wurden als wichtige Orientierungshilfe genannt, insbesondere fuer neue oder seltene Nutzende.

\textbf{Rollendifferenzierung} (3/10 TN) \\
Einzelne Teilnehmende wuenschten sich eine staerkere Sichtbarkeit rollenrelevanter Inhalte.

\textbf{Wahrgenommene Nutzerzentrierung} (7/10 TN) \\
UX-Strategien wurden implizit ueber Struktur, Konsistenz und Klarheit erkannt.

\textbf{Generische Gestaltung} (6/10 TN) \\
Eine rollenuebergreifende Grundstruktur wurde mehrheitlich akzeptiert.

% -----------------------------------------------------
\subsection{RQ-3: Akzeptanz und organisationale Einbettung}

\textbf{Akzeptanztreiber} \\
Aktualitaet der Inhalte (6 TN), zentrale Wissensquelle (5 TN), klare Struktur (5 TN), konkreter Arbeitsnutzen (4 TN).

\textbf{Akzeptanzbarrieren} \\
Veraltete Inhalte (6 TN), schwache Suchfunktion (5 TN), parallele Wissensquellen (4 TN), unklare Nutzungserwartung (3 TN).

\textbf{Kommunikationsbedarf} \\
Vorstellung in Meetings (8 TN), dialogische Formate (6 TN), regelmaessige Updates (6 TN), Fuehrungskraefte als Multiplikatoren (5 TN).

\textbf{Governance und Verantwortung} \\
Klare Zustaendigkeiten fuer Pflege und Qualitaetssicherung wurden mehrfach gefordert.

% -----------------------------------------------------
\subsection{Zusammenfassende Bewertung der Kategorien}

Die qualitative Analyse zeigt eine hohe Konsistenz zwischen subjektiver Wahrnehmung und objektiven Nutzungsdaten. Die Informationsarchitektur stellt den zentralen Hebel fuer Gebrauchstauglichkeit dar. Akzeptanz wird weniger durch Interface-Qualitaet als durch organisationale Rahmenbedingungen bestimmt. Die identifizierten Kategorien bilden die analytische Grundlage der Ergebnisdarstellung in Kapitel~5.
